%%═══════════════════════════════════════════════════════════════════
%% Lecture 14 — Stellar Structure as a BVP Application
%% 77315 — Computational Physics
%%═══════════════════════════════════════════════════════════════════

\section{Lecture 14 — Stellar Structure}\label{sec:lec14}
\hrule\vspace{1.5em}

%%──────────────────────────────────────────────────────────────────
\subsection*{Overview}
Stellar interiors are modelled by four coupled first-order ODEs
(the equations of stellar structure).  Because boundary conditions
are split between the centre and the surface, this is a BVP solved
by \emph{shooting to a fitting point} (Lecture~\ref{sec:lec13}).
We must find four unknown central values by requiring continuity of
all four quantities at the fitting point.

%%──────────────────────────────────────────────────────────────────
\subsection{The Four Stellar-Structure Equations}\label{subsec:lec14:odes}

Taking mass $m$ as the independent variable (instead of radius $r$):
\begin{align}
  \frac{dr}{dm} &= \frac{1}{4\pi r^2 \rho}, \label{eq:lec14:drm}\\
  \frac{dp}{dm} &= -\frac{G m}{4\pi r^4}, \label{eq:lec14:dpm}\\
  \frac{dL}{dm} &= \epsilon(\rho, T), \label{eq:lec14:dLm}\\
  \frac{dT}{dm} &= -\frac{3\kappa L}{64\pi^2 ac\,T^3 r^4}
                   \quad (\text{radiative transport}).
  \label{eq:lec14:dTm}
\end{align}

Here $\rho$ is density, $p$ pressure, $L$ luminosity, $T$ temperature,
$\epsilon$ energy generation rate per unit mass, $\kappa$ opacity,
$a$ radiation constant, $c$ speed of light.

%%──────────────────────────────────────────────────────────────────
\subsection{Boundary Conditions}\label{subsec:lec14:bc}

\subsubsection*{Centre ($m=0$)}
\begin{equation}
  r(0)=0,\quad L(0)=0.
  \label{eq:lec14:bc_centre}
\end{equation}
The free parameters are $\rho_c = \rho(0)$ and $T_c = T(0)$.

\subsubsection*{Surface ($m=M$)}
\begin{equation}
  p(M)=0,\quad T(M) = T_{\rm eff}.
  \label{eq:lec14:bc_surface}
\end{equation}
Two conditions at the surface, two unknown central values
$\Rightarrow$ a 2-parameter shooting problem.

%%──────────────────────────────────────────────────────────────────
\subsection{Shooting to a Fitting Point}\label{subsec:lec14:fitting}

\begin{enumerate}
  \item \textbf{Inner integration}: integrate outward from $m=0^+$
        to $m=m_f$ (fitting point, typically $m_f=M/2$) using
        guesses $\rho_c$, $T_c$.
  \item \textbf{Outer integration}: integrate inward from $m=M$
        to $m=m_f$ using surface conditions.
  \item \textbf{Matching}: require $r$, $p$, $L$, $T$ to be
        continuous at $m_f$.  This gives 4 equations for the
        4 unknowns $\rho_c$, $T_c$, $p_s$, $T_s$ (or parameters
        absorbed by the surface conditions).
  \item Solve the matching system with multidimensional Newton–Raphson.
\end{enumerate}

\nt{The shooting-from-both-ends approach avoids the exponential growth of solutions that makes one-sided shooting numerically unstable for stellar structure.}

%%──────────────────────────────────────────────────────────────────
\subsection{Energy Transport: Radiative vs.\ Convective}\label{subsec:lec14:transport}

Equation~\eqref{eq:lec14:dTm} assumes radiative energy transport.
When the actual temperature gradient
$\nabla_\text{rad} = (d\ln T/d\ln P)_\text{rad}$ exceeds the
adiabatic gradient $\nabla_\text{ad} = (\gamma-1)/\gamma$,
the layer is \emph{convectively unstable} (Schwarzschild criterion).
In convective zones, replace $dT/dm$ by the adiabatic gradient:
\begin{equation}
  \frac{dT}{dm}\bigg|_\text{conv}
  = \frac{T}{p}\nabla_\text{ad}\frac{dp}{dm}.
  \label{eq:lec14:conv}
\end{equation}
Check the stability criterion at each point during the integration.

%%──────────────────────────────────────────────────────────────────
\subsection{Chemical Evolution (Overview)}\label{subsec:lec14:chem}

Over stellar lifetimes, nuclear burning changes the composition.
A fifth equation tracks the hydrogen mass fraction $X$:
\begin{equation}
  \frac{dX}{dt} = -\frac{\epsilon_{pp}\rho}{Q_{pp}},
  \label{eq:lec14:chem}
\end{equation}
where $Q_{pp}$ is the energy released per proton–proton reaction.
This couples the structure equations to a time evolution,
turning the problem into a sequence of quasi-static BVPs.
